% from ega3-1-fr.tex
\oldpage[0\textsubscript{III}]{64}
flasques (G, II, 5.2.3). On a par suite
(\hyperref[0.11.1.6-fr]{11.1.6}) un isomorphisme canonique
\[
H^n(\check C^\bullet(\mathfrak U,\mathcal L^{\bullet\bullet}))
\xrightarrow{\sim}
H^n(\Gamma(X,\mathcal L^{\bullet\bullet}))
\]
(en vertu de (G, II, 5.2.2)), donc par définition
(\hyperref[0.12.4.2-fr]{12.4.2}) un isomorphisme
\[
H^n(\check C^\bullet(\mathfrak U,\mathcal L^{\bullet\bullet}))
\xrightarrow{\sim}\mathbf H^n(X,\mathcal K^\bullet).
\]
Considérons d'autre part le tricomplexe
$\check C^\bullet(\mathfrak U,\mathcal L^{\bullet\bullet})$ comme un
bicomplexe pour les degrés $i+j$ et $k$. Comme $k$ ne prend que des valeurs
$\geq0$, la première suite spectrale de ce bicomplexe est toujours
régulière~; elle est birégulière si $\mathcal L^{jk}=0$ pour $j<j_0$,
c'est-à-dire lorsque $\mathcal K^\bullet$ est limité inférieurement
(\hyperref[0.11.3.3-fr]{11.3.3}). Cette suite spectrale est la suite
cherchée~; en effet, pour tout $j$, $\mathcal L^{j,\bullet}$ est une
résolution injective de $\mathcal K^j$~; par suite,
$H^q(\check C^i(\mathfrak U,\mathcal L^{j,\bullet}))$ n'est autre que le
complexe de cochaînes
$\check C^i(\mathfrak U,h^q(\mathcal K^j))$, ce qui termine la
démonstration.
\end{proof}

\begin{corollary}[12.4.7]
\label{0.12.4.7-fr}
\emph{Si, pour tout simplexe $\sigma$ du nerf de $\mathfrak U$, et pour
tout entier $i$, on a
$H^q(U_\sigma,\mathcal K^i)=0$ pour $q>0$, alors on a un isomorphisme
canonique}
\begin{equation}
\label{0.12.4.7.1-fr}
\tag{12.4.7.1}
\mathbf H^\bullet(\mathfrak U,\mathcal K^\bullet)
\xrightarrow{\sim}\mathbf H^\bullet(X,\mathcal K^\bullet).
\end{equation}
\end{corollary}

\begin{proof}
En effet, l'hypothèse entraîne que
$\check C^\bullet(\mathfrak U,h^q(\mathcal K^i))=0$ pour $q>0$, donc
$E_2^{pq}=0$ pour $q>0$~; la suite
(\hyperref[0.12.4.6.1-fr]{12.4.6.1}) étant dégénérée et régulière, la
conclusion résulte de la définition
(\hyperref[0.12.4.5-fr]{12.4.5}) de
$\mathbf H^\bullet(\mathfrak U,\mathcal K^\bullet)$ et de
(\hyperref[0.11.1.6-fr]{11.1.6}).
\end{proof}

\begin{env}[12.4.8]
\label{0.12.4.8-fr}
Soit $(X',\mathcal O_{X'})$ un second espace annelé, et soit
$f=(\psi,\theta)$ un morphisme de $X'$ dans $X$. Par la même méthode que
dans (\hyperref[0.12.1.3-fr]{12.1.3}) et
(\hyperref[0.12.1.4-fr]{12.1.4}), on définit un di-homomorphisme pour
l'hypercohomologie d'un complexe $\mathcal K^\bullet$ de
$\mathcal O_X$-Modules
\begin{equation}
\label{0.12.4.8.1-fr}
\tag{12.4.8.1}
\mathbf H^p(X,\mathcal K^\bullet)
\longrightarrow\mathbf H^p(X',f^*(\mathcal K^\bullet)).
\end{equation}
On part d'une résolution injective de Cartan-Eilenberg
$\mathcal L^{\bullet\bullet}$ de $\mathcal K^\bullet$, et comme $\psi^*$
est exact, $\psi^*(\mathcal L^{\bullet\bullet})$ est une résolution de
Cartan-Eilenberg de $\psi^*(\mathcal K^\bullet)$ dans la catégorie des
$\psi^*(\mathcal O_X)$-Modules~; il y a alors un morphisme
$\psi^*(\mathcal L^{\bullet\bullet})\to
\mathcal L'^{\bullet\bullet}$, où $\mathcal L'^{\bullet\bullet}$ est une
résolution injective de Cartan-Eilenberg de
$\psi^*(\mathcal K^\bullet)$, et on en déduit un morphisme pour la
cohomologie~:
\[
\mathbf H^\bullet(X,\mathcal K^\bullet)
\longrightarrow\mathbf H^\bullet(X',\psi^*(\mathcal K^\bullet))~;
\]
par composition avec le morphisme déduit par fonctorialité de
$\psi^*(\mathcal K^\bullet)\to f^*(\mathcal K^\bullet)$, on obtient le
morphisme (\hyperref[0.12.4.8.1-fr]{12.4.8.1}) cherché.

Partant de (\hyperref[0.12.4.8.1-fr]{12.4.8.1}) et de
(\hyperref[0.12.4.3-fr]{12.4.3}), on peut alors, en raisonnant comme dans
(\hyperref[0.12.2.5-fr]{12.2.5}), définir, pour deux morphismes
$f:X\to Y$, $g:Y\to Z$ d'espaces annelés, des homomorphismes pour
l'hypercohomologie d'un complexe $\mathcal K^\bullet$ de
$\mathcal O_X$-Modules
\begin{equation}
\label{0.12.4.8.2-fr}
\tag{12.4.8.2}
\mathcal H^n(g,f_*(\mathcal K^\bullet))
\longrightarrow\mathcal H^n(h,\mathcal K^\bullet)
\end{equation}
\begin{equation}
\label{0.12.4.8.3-fr}
\tag{12.4.8.3}
\mathcal H^n(h,\mathcal K^\bullet)
\longrightarrow g_*(\mathcal H^n(f,\mathcal K^\bullet)).
\end{equation}
Nous laissons au lecteur le détail des définitions.
\end{env}

\section{Limites projectives en algèbre homologique}
\label{section:0.13-fr}

\subsection{La condition de Mittag-Leffler}
\label{subsection:0.13.1-fr}

\begin{env}[13.1.1]
\label{0.13.1.1-fr}
Soit $\mathcal C$ une catégorie abélienne dans laquelle les produits infinis
existent (axiome AB~3* de T, 1.5))~; alors la borne inférieure d'une famille
de sous-objets d'un objet de $\mathcal C$ existe, et tout système projectif
d'objets de $\mathcal C$ admet une limite projective, qui est un foncteur
exact à gauche du système projectif considéré (T, 1.8). Soit
$(A_\alpha,f_{\alpha\beta})$
